\paragraph{Statement A.4.5 (Proved endpoint-strip rigidity).}
Put $b=b_0$, $A=A_z=\sqrt{9z^2-1}$ and
$\delta=\operatorname{arsinh}(1/A)$. For an integral solution
$s^2-3t^2=2A^2$, $s>0$, write
$s=\sqrt2A\cosh u$, $t=\sqrt{2/3}A\sinh u$.
After choosing the lift of $u$ nearest the indicated endpoint,
\[
 |u|\leq\delta\ \Longrightarrow\ t=0,
 \qquad |u-b|\leq\delta\ \Longrightarrow\
 \bigl(4s-17z=0\text{ or }(s,t,z)=(4,0,1)\bigr).
\]
Consequently only $O(\log(2X))$ reduced ideals with $z\leq X$
can lie within these two endpoint strips. These are exactly the
possible endpoint-crossing terms in the spherical correction of
Statement A.4.4; this assertion concerns the geometric support of
that correction, not its complete amplitude.
\begin{proof}
The first implication follows from
$|t|\leq\sqrt{2/3}A\sinh\delta=\sqrt{2/3}<1$.
For the second use
\[
 \cosh b=\frac{17}{12\sqrt2},\qquad
 \sinh b=\frac1{12\sqrt2},\qquad
 \cosh\delta=\frac{3z}{A}.
\]
It follows that
\[
 \sqrt2A\cosh(b\pm\delta)=\frac{17z}{4}\pm\frac1{12}.
\]
Since $b>0$ and $\cosh$ is increasing on $[0,\infty)$, this
bounds $s$ when $\delta<b$, since the entire strip is positive.
In the remaining case $\delta\geq b$, we have $A\leq
1/\sinh b=12\sqrt2$, hence $z\leq5$. Directly for these five
heights, the only integral solutions in the second strip are
$(4,0,1)$ and $(17,1,4)$; equivalently one can note that the negative part lies
in the first strip and has $t=0$, and test $s^2=18z^2-2$ for
$z=1,2,3,4,5$. The only $t=0$ point in this finite range is $(4,0,1)$.
Outside this one stated exception, $|4s-17z|\leq1/3$, and integrality forces
$4s=17z$. The latter equation gives $z=4a$, $s=17a$ and
$a^2-3t^2=-2$, precisely the $n=0$ Pell family already counted.
The first equation gives $s^2-18z^2=-2$, another Pell equation;
its positive solutions form finitely many unit orbits, hence have
$O(\log(2X))$ members up to height $X$. Alternatively, ordinary
reduction by the fixed unit $17+12\sqrt2$ proves this bound directly.

Finally, in \eqref{A:F-spherical-integral} the multiplier of
$e^{i\kappa u}$ is an integral of
$e^{-i\kappa\log B(x)/2}$, with
$B(x)=\cosh(2\delta)+x\sinh(2\delta)$ and $-1\leq x\leq1$.
Since $e^{-2\delta}\leq B(x)\leq e^{2\delta}$, it only shifts
$u$ by a number in $[-\delta,\delta]$. An interval indicator
can therefore change only in the strips just considered. There
is also the scalar factor $B(x)^{-s/2}$: away from the strips
it gives $F_{s,0}(A)$, rather than $1$. No disappearance of that
amplitude correction is claimed.
\end{proof}

\paragraph{Statement A.4.6 (Proved existence of the sharp spectral arc functional).}
For each fixed nonconstant eigenfunction $u_j$, the limit
\[
 Q_{j,I}=\lim_{H\to\infty}
 \sum_{|\ell|\leq H}\widehat{\mathbf1_I}(\ell)
 \frac{D_{s_j}(0)}{D_{s_j}(\kappa_\ell)}p_{j,\ell}
\]
exists. Here $I=[0,b_0/L_0]$, the Fourier coefficient is with
respect to $\theta=u/L_0$, and the periods are unnormalized as
in Statement A.4.4. The limit is a well-defined interval value
of a distribution whose local primitive is H\"older continuous
with exponent $1/2$. This proves existence, not nonvanishing,
of the actual arc coefficient.
\begin{proof}
Conjugate the lift of the source axis to $\{ie^u:u\in\mathbb R\}$
in the upper half-plane. The Poisson--Helgason boundary
representation writes the lifted eigenfunction as
\[
 u_j(z)=\int_{\mathbb R}
       \left(\frac{\operatorname{Im}z}{|z-x|^2}\right)^{s_j}
       T_j(dx).
\]
This is a distributional integral, with the usual additional
chart at infinity. The precise theorem used is the boundary-value
representation in Anantharaman--Zelditch
\cite[Theorem 3.2]{AnantharamanZelditch2012}, together with their
Theorem 7.1 (Otal's boundary regularity theorem): a bounded
hyperbolic Laplace eigenfunction with bounded gradient and
parameter $s$, $\operatorname{Re}s\geq0$, has boundary values
which are the distributional derivative of a
$\operatorname{Re}s$-H\"older function. Our eigenfunction and
its gradient are bounded because the quotient is compact;
Statement A.2.5 gives $\operatorname{Re}s_j=1/2$. The source is
written for compact torsion-free quotients, so apply it to the
fixed cover of Statement A.3.1. The lifted eigenfunction is
unchanged as a function on the universal cover.

For $x=\pm e^v$, define the sum of two boundary distributions
\[
 \beta_j(dv)=(2e^v)^{-s_j}
       \bigl(T_j(d(e^v))+T_j(d(-e^v))\bigr),
\]
where $T_j(d(-e^v))$ denotes pushforward under
$x\mapsto\log|x|$ on the negative half-line, without an
orientation sign. Smooth changes of variable and multiplication
by a smooth function preserve the property of being the
derivative of a locally $1/2$-H\"older function. To see the
multiplication assertion explicitly, if $T=F'$ then
$aT=(aF)'-a'F$, and the second term has a continuously
differentiable primitive. The same formula after a smooth
coordinate change proves the other assertion.

The dilation $x\mapsto e^{L_0}x$ belongs to the source
stabilizer. Uniqueness of the Poisson boundary representation
implies $T_j(e^{L_0}dx)=e^{L_0s_j}T_j(dx)$, so $\beta_j$ is
$L_0$-periodic. Since
\[
 \left(\frac{e^u}{e^{2v}+e^{2u}}\right)^{s_j}
     =(2e^v)^{-s_j}\cosh(u-v)^{-s_j},
\]
its convolution with $\cosh(u)^{-s_j}$ is the restriction
$u_j(ie^u)$. The convolution is legitimate for the periodic
distribution because this kernel and each of its derivatives
decay exponentially. Taking its Fourier coefficients gives
\[
 p_{j,\ell}=K_{s_j}(\kappa_\ell)\,\beta_{j,\ell},
 \qquad K_s(\kappa)=2^{s-1}D_s(\kappa),\qquad
 \beta_{j,\ell}=\int_0^{L_0}e^{i\kappa_\ell v}\beta_j(dv).
\]
Therefore the displayed partial sum defining $Q_{j,I}$ equals
$K_{s_j}(0)$ times the pairing of $\beta_j$ with the symmetric
Fourier partial sum of $\mathbf1_{[0,b_0]}$.

Let $F_j$ be a continuous local primitive of $\beta_j$;
subtracting its affine mean gives a periodic $1/2$-H\"older
function. Its Fourier partial sums converge uniformly. For
completeness, convolving a H\"older function with the Fej\'er
kernel gives uniform error $O_j(H^{-1/2})$ by splitting the
kernel integral at $1/H$; composing the approximation with the
Dirichlet projection of degree $H$ and using its elementary
$L^1$ bound $O(\log(2H))$ gives an error
$O_j(H^{-1/2}\log(2H))$ for the ordinary partial sum (use a
Fej\'er polynomial of degree at most $H$). Integration by parts
then identifies the limit with
\[
 Q_{j,I}=K_{s_j}(0)\bigl(F_j(b_0)-F_j(0)\bigr).
\]
This value is independent of the primitive's additive constant,
and the convergence error is $O_j(H^{-1/2}\log(2H))$.
The implied constant here depends on the fixed eigenfunction;
no exchange with an unbounded spectral sum follows from this
fixed-eigenfunction assertion alone.
\end{proof}
